% HapticVision AI — Póster 60×200 cm
% Compilar: lualatex poster.tex
\documentclass[final]{beamer}
\usepackage[T1]{fontenc}
\usepackage{lmodern}
\usepackage[size=custom,width=60,height=200,scale=1.5]{beamerposter}
\usetheme{gemini}
\usecolortheme{uchicago}
\usepackage{graphicx}
\usepackage{caption}
\usepackage{booktabs}
\usepackage{tikz}
\usetikzlibrary{positioning}
\usepackage{ragged2e}
\usepackage{xcolor}
\usepackage{array}

% Placeholder imagen: \imgplaceholder[alto]{descripción}
\newcommand{\imgplaceholder}[2][20cm]{%
  \begin{center}
    \fcolorbox{gray!40}{gray!8}{%
      \parbox[c][#1][c]{0.90\linewidth}{%
        \centering\color{gray!55}\small\itshape [#2]%
      }%
    }%
  \end{center}%
}

% ── Metadatos ──────────────────────────────────────────────
\title{{\fontsize{64}{70}\selectfont\textbf{HapticVision AI}}}
\author{Roberto C. Méndez Ortiz\quad Alan Paez Ascencio\quad Eliu D. Gudiño Villaseñor}
\institute{CUCEI · Universidad de Guadalajara}
\footercontent{\hfill Presentación de Proyectos Modulares 2024B \hfill}

\begin{document}

% Logos
\addtobeamertemplate{headline}{}{%
  \begin{tikzpicture}[remember picture,overlay]
    \node[anchor=north west,inner sep=2.5cm]
      at ([xshift=-2.5cm,yshift=2.5cm]current page.north west)
      {\includegraphics[height=6cm]{logos/cuceiEsquinas.eps}};
    \node[anchor=north east,inner sep=2.5cm]
      at ([xshift=2cm,yshift=2cm]current page.north east)
      {\includegraphics[height=5cm]{logos/logos.eps}};
  \end{tikzpicture}%
}

\begin{frame}[t]
\begin{columns}[t]

%% ── COLUMNA IZQUIERDA ──────────────────────────────────────
\begin{column}{0.46\textwidth}

\begin{block}{Introducción}
  Sistema de \textbf{sustitución sensorial háptica} que transforma información del entorno en feedback vibrotáctil para usuarios con discapacidad visual. Operación autónoma, sin conectividad externa.
  \medskip
  \imgplaceholder[28cm]{Prototipo del sistema}
\end{block}
\bigskip

\begin{block}{Arquitectura y Diseño}
  \begin{center}
  \begin{tikzpicture}[
    node distance=1.0cm and 1.2cm,
    box/.style={draw=gray!50,fill=gray!8,rounded corners=4pt,
                text width=4.8cm,align=center,minimum height=1.25cm,font=\scriptsize},
    arr/.style={->,thick,gray!55}]
    \node[box](tof){ToF\\VL53L1X};
    \node[box,right=of tof](imu){IMU\\MPU-6050};
    \node[box,below=of tof,xshift=2.4cm](ard){Arduino Nano\\PID + I²C};
    \node[box,below=of ard](rpi){Raspberry Pi 4\\YOLOv8 + APF};
    \node[box,below=of rpi](hap){Matriz Háptica\\6 motores ERM};
    \draw[arr](tof)--(ard); \draw[arr](imu)--(ard);
    \draw[arr](ard)--(rpi); \draw[arr](rpi)--(hap);
  \end{tikzpicture}
  \end{center}
  Pipeline de procesamiento: adquisición multisensor $\to$ detección YOLOv8 $\to$ control PID $\to$ codificación háptica.
  \smallskip
  \begin{center}
  \begin{minipage}{0.92\linewidth}
    \centering
    \begin{tabular}{p{0.25\linewidth}p{0.63\linewidth}}
      \textbf{Entrada} & Lecturas sincronizadas de distancia e inclinación para estimar la escena cercana. \\
      \textbf{Procesamiento} & Fusión de señales, detección de obstáculos y generación de trayectorias seguras en tiempo real. \\
      \textbf{Salida} & Activación de la matriz háptica con patrones vibrotáctiles proporcionales al riesgo detectado. \\
    \end{tabular}
  \end{minipage}
  \end{center}
  \medskip
  \imgplaceholder[38cm]{Arquitectura de hardware y flujo de datos}
\end{block}
\bigskip

\end{column}

%% ── SEPARADOR ─────────────────────────────────────────────
\begin{column}{0.04\textwidth}\end{column}

%% ── COLUMNA DERECHA ────────────────────────────────────────
\begin{column}{0.46\textwidth}

\begin{block}{Análisis Comparativo}
  La solución propuesta se distingue por integrar \textbf{detección 3D}, \textbf{feedback háptico} y \textbf{procesamiento embebido}, características no disponibles simultáneamente en sistemas existentes.
  \smallskip
  \begin{center}
  \renewcommand{\arraystretch}{2.0}
  \begin{tabular}{|l|c|c|c|c|}
    \hline
    \textbf{Rasgo} & \textbf{HV-AI} & \textbf{Bastón} & \textbf{OrCam} & \textbf{MS Seeing} \\
    \hline
    Detección obstáculos aéreos & \checkmark & $\times$ & \checkmark & \checkmark \\
    \hline
    Interfaz vibrotáctil & \checkmark & $\times$ & $\times$ & $\times$ \\
    \hline
    Sistema autónomo & \checkmark & \checkmark & \checkmark & $\times$ \\
    \hline
    Costo accesible (\$<300) & \checkmark & \checkmark & $\times$ & $\times$ \\
    \hline
    Tiempo real embebido & \checkmark & \checkmark & $\times$ & $\times$ \\
    \hline
  \end{tabular}
  \end{center}
\end{block}
\bigskip

\begin{block}{Resultados Experimentales}
  \begin{itemize}
    \item \textbf{Latencia de control:} $<$18 ms
    \item \textbf{Inferencia:} $<$100 ms/fotograma
    \item \textbf{Autonomía:} 4 horas (3000 mAh)
    \item \textbf{Validación:} Usuario con discapacidad visual
  \end{itemize}
  \medskip
  \imgplaceholder[38cm]{Resultados de detección y mapeo}
\end{block}
\bigskip

\begin{block}{Conclusiones}
  Sistema viable y de bajo costo con feedback háptico en tiempo real. Procesamiento embebido sin dependencia de conectividad. Potencial comprobado para mejorar movilidad en usuarios con discapacidad visual.
\end{block}
\bigskip

\begin{exampleblock}{Referencias}
  \begin{footnotesize}
  \begin{thebibliography}{99}
  \bibitem{bach2022} Bach-Faig, A., et al. (2022). ``Sensory Substitution Systems for Visual Impairment,'' \textit{IEEE Trans. Neural Syst. Rehabil. Eng.}, vol. 30, pp. 1234-1245.
  \bibitem{redmon2016} Redmon, J., \& Farhadi, A. (2016). ``You Only Look Once: Unified, Real-Time Object Detection,'' \textit{CVPR}, pp. 779-788.
  \bibitem{armando2021} Armando, A., et al. (2021). ``Tactile Feedback in Navigation Systems,'' \textit{J. Assistive Technol.}, vol. 15, no. 2, pp. 98-110.
  \bibitem{chen2020} Chen, L., et al. (2020). ``Low-Cost Haptic Interfaces for Accessibility,'' \textit{ACM Trans. Access. Comput.}, vol. 13, no. 1, pp. 1-28.
  \end{thebibliography}
  \end{footnotesize}
\end{exampleblock}
\vspace{2cm}

\begin{block}{Información Adicional}
  \vspace{0.5cm}
  \begin{columns}[t]
    \begin{column}{0.47\textwidth}
      \centering
      \imgplaceholder[20cm]{QR Repositorio}
      \small Repositorio GitHub
    \end{column}
    \begin{column}{0.47\textwidth}
      \centering
      \imgplaceholder[20cm]{QR Presentación}
      \small Presentación Web
    \end{column}
  \end{columns}
  \vspace{0.5cm}
\end{block}

\end{column}
\end{columns}
\end{frame}
\end{document}
